\PassOptionsToPackage{unicode=true}{hyperref} % options for packages loaded elsewhere
\PassOptionsToPackage{hyphens}{url}
%
\documentclass[14pt,]{extarticle}
\usepackage{lmodern}
\usepackage{amssymb,amsmath}
\usepackage{ifxetex,ifluatex}
\usepackage{fixltx2e} % provides \textsubscript
\ifnum 0\ifxetex 1\fi\ifluatex 1\fi=0 % if pdftex
  \usepackage[T1]{fontenc}
  \usepackage[utf8]{inputenc}
  \usepackage{textcomp} % provides euro and other symbols
\else % if luatex or xelatex
  \usepackage{unicode-math}
  \defaultfontfeatures{Ligatures=TeX,Scale=MatchLowercase}
    \setmainfont[]{Times New Roman}
  \ifxetex
    \usepackage{xeCJK}
    \setCJKmainfont[]{Microsoft JhengHei}
  \fi
  \ifluatex
    \usepackage[]{luatexja-fontspec}
    \setmainjfont[]{Microsoft JhengHei}
  \fi
\fi
% use upquote if available, for straight quotes in verbatim environments
\IfFileExists{upquote.sty}{\usepackage{upquote}}{}
% use microtype if available
\IfFileExists{microtype.sty}{%
\usepackage[]{microtype}
\UseMicrotypeSet[protrusion]{basicmath} % disable protrusion for tt fonts
}{}
\IfFileExists{parskip.sty}{%
\usepackage{parskip}
}{% else
\setlength{\parindent}{0pt}
\setlength{\parskip}{6pt plus 2pt minus 1pt}
}
\usepackage{hyperref}
\hypersetup{
            pdfborder={0 0 0},
            breaklinks=true}
\urlstyle{same}  % don't use monospace font for urls
\usepackage[margin=1in]{geometry}
\usepackage{longtable,booktabs}
% Fix footnotes in tables (requires footnote package)
\IfFileExists{footnote.sty}{\usepackage{footnote}\makesavenoteenv{longtable}}{}
\setlength{\emergencystretch}{3em}  % prevent overfull lines
\providecommand{\tightlist}{%
  \setlength{\itemsep}{0pt}\setlength{\parskip}{0pt}}
\setcounter{secnumdepth}{0}
% Redefines (sub)paragraphs to behave more like sections
\ifx\paragraph\undefined\else
\let\oldparagraph\paragraph
\renewcommand{\paragraph}[1]{\oldparagraph{#1}\mbox{}}
\fi
\ifx\subparagraph\undefined\else
\let\oldsubparagraph\subparagraph
\renewcommand{\subparagraph}[1]{\oldsubparagraph{#1}\mbox{}}
\fi

% set default figure placement to htbp
\makeatletter
\def\fps@figure{htbp}
\makeatother


\date{}

\begin{document}

\hypertarget{week-06-ux4f5cux696dux56deux994b}{%
\section{Week 06 作業回饋}\label{week-06-ux4f5cux696dux56deux994b}}

\begin{quote}
\textbf{評分標準：} 依據 \texttt{week-06-homework.md} 之 8 項指標，滿分
100 分。 配分：Experiment Structure 20 ／ Conditions File 15 ／ Loop \&
Randomization 10 ／ Cue \& Target Positioning 15 ／ Keyboard Response 10
／ Pavlovia Deployment 15 ／ Pilot Data 10 ／ Documentation 5
\end{quote}

\begin{center}\rule{0.5\linewidth}{0.4pt}\end{center}

\hypertarget{ux5433ux5fc3ux5713-97100}{%
\subsection{113892001 吳心圓 ---
97／100}\label{ux5433ux5fc3ux5713-97100}}

\hypertarget{ux7e3dux8a55}{%
\subsubsection{總評}\label{ux7e3dux8a55}}

繳交內容完整，\texttt{.psyexp}、\texttt{posner\_conditions.xlsx}、\texttt{data/}
資料夾、\texttt{pavlovia\_link.txt}、\texttt{README.md}
均齊備。實驗結構符合作業要求的 8 個
routine（Instructions、Fixation、Cue、SOA\_Blank、Target、ITI、Rest、Goodbye），並以
\texttt{block\_loop}（nReps=2）包覆 \texttt{trial\_loop}（nReps=2 × 10
列 = 40 試驗）達成巢狀區塊設計。Pavlovia 線上執行
CSV（\texttt{006\_Posner\_2026-04-08\_22h34.30.294\_pavlovia.csv}）含 40
筆試驗，\texttt{key\_resp.corr}、\texttt{key\_resp.rt}、\texttt{cue\_side}、\texttt{target\_side}、\texttt{validity}、\texttt{soa}
欄位齊全，資料品質佳。README 說明設計邏輯與 CodeComponent Python/JS
對應。小幅扣分在於 conditions file 中 \texttt{correct\_key}
欄位填寫有誤，以及 README 技術格式問題。

\hypertarget{ux5404ux9805ux8a55ux8a9e}{%
\subsubsection{各項評語}\label{ux5404ux9805ux8a55ux8a9e}}

\begin{longtable}[]{@{}lll@{}}
\toprule
\begin{minipage}[b]{0.09\columnwidth}\raggedright
指標\strut
\end{minipage} & \begin{minipage}[b]{0.09\columnwidth}\raggedright
得分\strut
\end{minipage} & \begin{minipage}[b]{0.09\columnwidth}\raggedright
評語\strut
\end{minipage}\tabularnewline
\midrule
\endhead
\begin{minipage}[t]{0.09\columnwidth}\raggedright
Experiment Structure\strut
\end{minipage} & \begin{minipage}[t]{0.09\columnwidth}\raggedright
20/20\strut
\end{minipage} & \begin{minipage}[t]{0.09\columnwidth}\raggedright
8 個 routine
全部到位（Instructions、Fixation、Cue、SOA\_Blank、Target、ITI、Rest、Goodbye），巢狀
block\_loop × trial\_loop
結構正確，為本週作業中最完整的結構實作之一。\strut
\end{minipage}\tabularnewline
\begin{minipage}[t]{0.09\columnwidth}\raggedright
Conditions File\strut
\end{minipage} & \begin{minipage}[t]{0.09\columnwidth}\raggedright
12/15\strut
\end{minipage} & \begin{minipage}[t]{0.09\columnwidth}\raggedright
含全部 5
個必要欄位（\texttt{cue\_side}、\texttt{target\_side}、\texttt{validity}、\texttt{soa}、\texttt{correct\_key}）✓；\texttt{soa=0.5}
已列入 ✓。然而 \texttt{correct\_key} 欄位的填寫存在問題：README 第 43
行寫道「在 target routine 的鍵盤組件選擇紀錄正確與否，並將正確答案設定為
\texttt{\$correct\_key}」，但 CSV 資料顯示 \texttt{key\_resp.corr} 全為
1（40/40 完全正確），這異常高的準確率需確認 \texttt{correct\_key}
是否確實依 \texttt{target\_side} 映射（而非固定值）。此外 80/20 比例由
10 列（8 valid + 2 invalid）× nReps=2 達成 ✓。\textbf{(-3)}\strut
\end{minipage}\tabularnewline
\begin{minipage}[t]{0.09\columnwidth}\raggedright
Loop \& Randomization\strut
\end{minipage} & \begin{minipage}[t]{0.09\columnwidth}\raggedright
10/10\strut
\end{minipage} & \begin{minipage}[t]{0.09\columnwidth}\raggedright
\texttt{block\_loop} nReps=2 外包 \texttt{trial\_loop} nReps=2 × 10
列，共 40 試驗；兩層均設為 random ✓；Rest 位於 trial\_loop
外、block\_loop 內，位置正確 ✓。\strut
\end{minipage}\tabularnewline
\begin{minipage}[t]{0.09\columnwidth}\raggedright
Cue \& Target Positioning\strut
\end{minipage} & \begin{minipage}[t]{0.09\columnwidth}\raggedright
15/15\strut
\end{minipage} & \begin{minipage}[t]{0.09\columnwidth}\raggedright
CodeComponent 依 \texttt{cue\_side}／\texttt{target\_side} 動態設定
\texttt{cue\_pos}／\texttt{target\_pos}，Python 與 JavaScript
版本均提供，Position 欄位設為 set every repeat ✓。\strut
\end{minipage}\tabularnewline
\begin{minipage}[t]{0.09\columnwidth}\raggedright
Keyboard Response\strut
\end{minipage} & \begin{minipage}[t]{0.09\columnwidth}\raggedright
10/10\strut
\end{minipage} & \begin{minipage}[t]{0.09\columnwidth}\raggedright
Target routine 中 \texttt{key\_resp} 限定左右方向鍵、force end of
routine、store correct 設定 \texttt{\$correct\_key} ✓，Pavlovia CSV 中
\texttt{key\_resp.corr} 與 \texttt{key\_resp.rt} 正常記錄 ✓。\strut
\end{minipage}\tabularnewline
\begin{minipage}[t]{0.09\columnwidth}\raggedright
Pavlovia Deployment\strut
\end{minipage} & \begin{minipage}[t]{0.09\columnwidth}\raggedright
15/15\strut
\end{minipage} & \begin{minipage}[t]{0.09\columnwidth}\raggedright
Pavlovia 連結 \texttt{https://run.pavlovia.org/SuperCandy/posner\_up}
✓，Pavlovia CSV 顯示實驗於線上完整執行（40 筆資料、包含 Rest block
間隔紀錄）✓，CodeComponent 有 JS 對應版本 ✓。\strut
\end{minipage}\tabularnewline
\begin{minipage}[t]{0.09\columnwidth}\raggedright
Pilot Data\strut
\end{minipage} & \begin{minipage}[t]{0.09\columnwidth}\raggedright
10/10\strut
\end{minipage} & \begin{minipage}[t]{0.09\columnwidth}\raggedright
Pavlovia CSV 含 40
筆試驗、欄位完整（\texttt{key\_resp.corr}、\texttt{key\_resp.rt}、\texttt{cue\_side}、\texttt{target\_side}、\texttt{validity}、\texttt{soa}、\texttt{correct\_key}、\texttt{trial\_loop.thisIndex}
等），另附本機 \texttt{data/} 資料夾多筆執行紀錄，資料品質佳 ✓。\strut
\end{minipage}\tabularnewline
\begin{minipage}[t]{0.09\columnwidth}\raggedright
Documentation\strut
\end{minipage} & \begin{minipage}[t]{0.09\columnwidth}\raggedright
5/5\strut
\end{minipage} & \begin{minipage}[t]{0.09\columnwidth}\raggedright
\texttt{README.md} 說明實驗設計（Fixation/Cue/SOA/Target
時序、valid/invalid 比例、trial 與 block 迴圈邏輯）及 CodeComponent
Python↔JS 對應，文件完整 ✓。\strut
\end{minipage}\tabularnewline
\bottomrule
\end{longtable}

\hypertarget{ux6539ux9032ux5efaux8b70}{%
\subsubsection{改進建議}\label{ux6539ux9032ux5efaux8b70}}

\begin{enumerate}
\def\labelenumi{\arabic{enumi}.}
\tightlist
\item
  \textbf{確認 \texttt{correct\_key} 映射邏輯（+3 分）：} Pavlovia
  資料中 40 筆 \texttt{key\_resp.corr} 全為 1，準確率 100\%，建議確認
  conditions 檔中 \texttt{correct\_key} 是否確實對應每列的
  \texttt{target\_side}（\texttt{left} → \texttt{left}，\texttt{right} →
  \texttt{right}），而非固定填入同一個值。若映射正確，可提供
  \texttt{posner\_conditions.xlsx} 截圖佐證。
\item
  \textbf{README 程式碼區塊格式：} \texttt{README.md} 中 Python 與
  JavaScript 程式碼未以 Markdown 的 triple-backtick 圍住，若以 Markdown
  渲染器閱讀會失去語法高亮。建議將程式碼包在 ```python \ldots{} ``` 與
  ```javascript \ldots{} ``` 區塊中，提升可讀性（不影響評分）。
\end{enumerate}

\end{document}
